\section*{对象、模型与结构。}
\phantomsection
\addcontentsline{toc}{section}{对象、模型与结构。}

\begin{enumerate}
\def\labelenumi{\Alph{enumi})}
\tightlist
\item
  数学的对象与结构。——从古代直到 19 世纪，关于数学家的主要对象是什么，存在着一种共同的理解；它们恰恰就是柏拉图在前文所引段落（第~13~页）中提到的那些：数、量与图形。如果说起初还必须加上作为力学、天文学、光学和音乐之题材的那些对象和现象，那么在希腊人那里，这些“数学性”的学科也总是与算术和几何清楚地分开，而从文艺复兴开始，它们相当快地取得了独立科学的地位。
\end{enumerate}

无论数学对象的概念被这样那样的数学家或哲学家的哲学细微色彩如何渲染，至少在一点上是一致的：这些对象是给定的，我们无权指派给它们任意的性质，正如物理学家不能改变一种自然现象一样。说实在的，心理反应无疑构成了这些看法的一部分，这些反应不是我们要追踪的，但每个数学家在徒劳地力图抓住一个似乎不断溜走的证明而耗尽心力时，都体会过它们。从这里到把这种阻力同感官世界横亘在我们路上的障碍等同起来，只有一步之遥；即使在今天，不止一位自称不妥协的形式主义者，在其内心深处也会乐于赞同埃尔米特的这段自白：“我相信，数和解析函数并不是我们心智的任意产物；我认为它们存在于我们之外，带有与客观实在的事物同样的必然性品格；我们遇见它们或发现它们，研究它们，就像物理学家、化学家和动物学家那样”（{[}160{]}，第 II 卷，第~398~页）。

在数学的经典观念里，根本谈不上游离于对数与图形的研究；但是，这个每一位数学家都认为自己必须在口头上表示赞同的官方学说，随着新观念的积累，却渐渐变成一种无法忍受的负担。代数学家面对负数的窘迫，直到解析几何给它提供了一个方便的“解释”才告结束；然而迟至 18 世纪，达朗贝尔在《百科全书》中（{[}75 a{]}，“NÉGATIF”条）讨论这个问题时，在一大段相当含混的解释之后突然没了底气，只满足于得出这样的结论：“关于负量的代数运算规则为人们普遍假定，并被普遍视为正确，无论人们对这些量另有何种想法。”至于虚数，丑闻要大得多；因为，如果它们是“不可能的”根，而且（直到 1800 年前后）看不到任何“解释”它们的途径，那么人们怎么能不出矛盾地谈论这些无法定义的存在，尤其是为什么要引入它们呢？达朗贝尔在这里保持着谨慎的沉默，甚至不曾提出这些问题，无疑是因为他意识到，自己无法作出比一个世纪前 A. 吉拉尔的天真回答更好的回答（{[}129{]}，f.~22）：“也许有人会问：这些不可能的解有什么用？我回答：有三点用处：为了保证一般规则，为了表明没有其他的解，也为了它本身的用处。”

在分析方面，17 世纪的局面也不见得好。一个幸运的情况是，解析几何仿佛应约而至地出现，以几何图形的形状为 17 世纪的伟大创造——函数概念——提供了一个“表示”，从而（与费马、帕斯卡和巴罗一道）有力地促成了无穷小演算的诞生（参见第~193~页）。但另一方面，人们也知道无穷小和不可分量这些观念将引起何等哲学-数学的论争。如果说达朗贝尔在这里要幸运些——他承认无穷小演算的“形而上学”中除了极限概念别无他物（{[}75 a{]}，“DIFFÉRENTIEL”条和“LIMITE”条，以及 {[}75 b{]}）——那么他仍不能比他的同时代人更明白发散级数展开的真实含义，也不能解释这样一个悖论：用剥除了一切数值解释的表达式算到最后，得到的却是精确的结果。最后，即使在“几何确定性”的领域内，欧几里得的框架也爆裂了：当斯特林在 1717 年不惮于说某条曲线有一个“二重虚无穷远点”时（{[}299{]}，新版第~93~页），要把这样一个“对象”与通常理解的概念联系起来，他肯定会有困难；而在 19 世纪初创立射影几何、使这类思想获得相当发展的彭赛列（见第~131~页），也仍然只满足于援引一个纯然形而上学的“连续性原理”来为自己辩护。

可以设想，在这些条件下（而且恰恰在“绝对真理”被以最强音宣布出来的那个时刻），证明的概念在 18 世纪中似乎变得越来越模糊，因为像希腊人那样，把进行推理所依据的概念及其基本性质一劳永逸地确定下来，已经无从谈起。始于 19 世纪初的向严格性的回归使这种状况有所改善，却没有因此挡住新观念的洪流：于是我们看到，代数中出现了伽罗瓦的虚数（{[}123{]}，第~113-127~页）、库默尔的理想数 {[}188 b{]}，继之而起的是向量和四元数、n 维空间、多重向量和张量（见第~61~页以下），更不消说布尔代数。无疑，重大的进步——正是它使严格性的回归得以不失去先前时代的任何收获——来自用更经典的术语为这些新观念给出“模型”的可能性：理想数或伽罗瓦虚数借助同余理论得到解释（见第~82~页以下）；n 维几何（如果愿意这样看）只表现为陈述“含 n 个变量”的代数结果的一种纯粹语言；至于经典的虚数——它以平面上点所作的几何表示（见第~161~页以下）标志着代数这场繁盛的开端——不久便可以在几何“模型”与用同余表述的解释之间任择其一（参见第~82~页）。但数学家们终于开始痛感，他们是在逆着一个自然的斜坡而奋斗，他们的工作正被这个斜坡拖着走；而在数学中，讨论没有任何合适“解释”的对象，应当是容许的：“数学的本质”，布尔在 1854 年说，“不在于埋头处理数与量的观念”（{[}29{]}，第 II 卷，第~13~页）。\(^{30}\) 同样的关切促使格拉斯曼在他的《Ausdehnungslehre》

1844 年版中，以一种起初排除了数或几何对象概念的形式来呈现他的演算。\(^{31}\) 而稍后，黎曼在他的就职演讲中，一开始就注意不说“点”，而在描述“n 重延展的流形”时说“确定方式”（Bestimmungsweise），并强调在这种流形中，“度量关系”（Massverhältnisse）“只能对抽象的量进行研究，只能用公式来表示；在一定的条件下，却可以把它们分解为一些关系，其中每一个孤立地看都容许一种几何表示，从而能够以几何的形式表达计算的结果”（{[}259 a{]}，第~276~页）。

从这时起，公理方法的扩张已是既成事实。如果说在一段时间里，人们还认为（只要可能）用几何直觉来检验“抽象”结果是有益的，那么至少已经公认，“经典”对象不再是数学家可以合法研究的唯一对象。这是因为——恰恰由于多种“解释”或“模型”皆属可能——人们已经认识到，数学对象的“本性”终究是次要的；比如，一个结果无论作为“纯”几何的定理呈现，还是借助解析几何作为代数定理呈现，都无关紧要。换句话说，数学的本质——这个直到那时只能用“通则”或“形而上学”之类含糊名称来表达的难以捉摸的概念——显现为对象之间关系的研究，而这些对象只是（有意地）通过其若干性质被了解和描述，恰好就是被当作公理置于其理论根基的那些性质。布尔在 1847 年就已经清楚地看到了这一点，他写道，数学处理的是“就其自身加以考虑的运算，独立于它们可以应用于其上的各种对象”（{[}29{]}，第 I 卷，第~3~页）。汉克尔在 1867 年开创代数的公理化时，捍卫的是一种“纯粹智性的数学，一种纯粹形式的理论，其目的不在于量、或量的影像即数的结合，而在于思想对象（“Gedankendinge”），可以有实在的对象或关系与它们对应，尽管这样的对应并非必要”（{[}146{]}，第~10~页）。康托尔在 1883 年以宣称“数学在其发展中是完全自由的，它的概念只受相容性这一必然性的约束，并通过精确定义与先前引入的概念相衔接”，来回应这一“自由数学”的主张（{[}47{]}，第~182~页）。最后，欧氏几何的修订成功地传播并普及了这些观念。帕施本人虽然仍执着于几何对象的某种“实在性”，也承认几何学实际上独立于这些对象的意义，纯粹在于研究它们的关系（{[}245{]}，第~90~页）；希尔伯特把这一观念推向其逻辑结论，强调甚至一个数学理论基本概念的名称也可以任意选择，\(^{32}\) 而庞加莱则把这一点表述为：公理是“伪装的定义”，从而完全颠倒了经院的观点。

因此，人们会禁不住说，现代的“结构”概念在实质上已于 1900 年前后达到；事实上，它还需要再经三十年左右的学徒期，才能以其全部光辉出现。当同类结构的性质相当简单时，辨认它们无疑并不困难；例如就群结构而言，这一阶段早在 19 世纪中叶便已达到。但与此同时，我们看到汉克尔在搏斗——却没有完全成功——以提炼出域和扩张的一般观念，他只能以一条半形而上学的“永久性原理”的形式来表达它 {[}146{]}，而这条原理要到 40 年后才由施泰尼茨 {[}294 a{]} 以最终的形式表述出来。尤其是，在这件事上，人们颇难摆脱这样一种印象：数学对象是连着其结构一起“给定”给我们的；只有对泛函分析相当长期的使用，才能使现代数学家熟悉如下想法：例如，有理数上有多个“自然的”拓扑，数直线上有多种测度。伴随着这种分离，向结构的一般定义的过渡终于实现了。

\begin{enumerate}
\def\labelenumi{\Alph{enumi})}
\setcounter{enumi}{1}
\item
  模型与同构。——用一门数学理论对另一门数学理论给出“模型”或“解释”这一观念的介入，我们已经多次注意到。这并不是一个新观念，而且我们在这里看到的，无疑是一种反复重现的感受的表现：诸门“数学科学”深层的统一性。如果可以把早期毕达哥拉斯学派的传统格言“万物皆数”视为可信，那么它可以被视为把当时的几何与代数归约为算术的第一次尝试的痕迹。尽管不可公度数的发现似乎永远关闭了这条道路，但它在希腊数学中引起的反应却是第二次综合的尝试，这一次以几何为基础，并把得自巴比伦人的代数方程解法等一并吸收进来。\(^{33}\) 众所周知，这一观念一直延续到 R. 邦贝利和笛卡尔的根本改革为止，后者把量的所有度量都同化为长度的度量（换句话说，同化为实数；参见第~151~页）。但随着笛卡尔和费马创立解析几何，倾向再次逆转，几何与代数得到了紧密得多的融合，不过这一次受益的是代数。此外，在别处，笛卡尔走得更远，构想出“所有这些通常被称为数学的科学 \ldots{} 尽管它们的题材各不相同”的本质统一性，“它们仍不会不汇合”，他说，“因为它们所考虑的，不过是其中可以找到的各种关系或比例”（{[}85 a{]}，第 VI 卷，第~19-20~页）。\(^{34}\) 然而这个观点只不过倾向于使代数成为基础的科学；莱布尼茨激烈地抗议这一结论，而他本人如前所述也构想了一种“普遍的数学”，但其规模宏大得多，而且已经相当接近现代观念。他把笛卡尔所说的“一致”精确化，实际上第一次窥见了同构的一般概念（他称之为“相似”），以及把同构的关系或运算“视为同一”的可能性；他给出了加法和乘法的例子（{[}69 a{]}，第~301-303~页）。但这些大胆的见解在他同时代人中没有引起任何回声，必须等到 19 世纪中叶前后代数的扩张（见第~51~页以下），才能看到莱布尼茨之梦开始实现。我们已经强调过，正是在那个时刻，“模型”成倍增加，通过简单的语言转换从一门理论过渡到另一门理论成为常事；其中最引人注目的例子也许是射影几何中的对偶（见第~132~页），当时把互为“对偶”的定理面对面排成两栏刊印的常见做法，无疑对同构概念的充分实现起了很大作用。从更技术的观点看，阿贝尔群的同构群概念当然已为高斯所知，置换群的同构群概念已为伽罗瓦所知（参见第~51~页以下）；对无论什么样的群，这一概念大约在 19 世纪中叶便已一般地获得。\(^{35}\) 此后，对每一门新的公理化理论，定义一个同构概念都是顺理成章的发展；但只是随着现代的结构概念，人们才最终认识到，每一结构都自带一个同构概念，无需对每种类型的结构分别给出专门的定义。
\item
  经典数学的算术化。——对“模型”观念越来越广泛的使用，还将使 19 世纪实现毕达哥拉斯学派所梦想的数学统一。在本世纪初，整数与连续量显得一如既往地不可调和，一如古代那样；实数仍然与几何量的概念（至少与长度的概念）联系在一起，而且人们正是求助于后者来构造负数和虚数的“模型”。甚至连有理数，也传统地附着于把一个量分成相等几份的想法；只有整数始终独立在外，正如高斯在 1832 年所说的那样，是“我们心智的专属产物”，他把它们与空间观念相对立（{[}124 a{]}，第 VIII 卷，第~201~页）。使算术与分析相互接近的最初努力，起初是围绕（正和负的）有理数进行的，出自马丁·欧姆（1822 年）；这些努力在 1860 年前后被几位作者重新提起，著名的有格拉斯曼、汉克尔和魏尔斯特拉斯（在其未刊的讲课中）；看来正是最后这一位想出了通过考虑整数对的类，来得到正有理数或负整数的“模型”的想法。但最重要的一步仍待迈出，即在有理数理论中为无理数找到一个“模型”；到 1870 年前后，这已成为一个紧迫的问题，因为在分析中发现“病态”现象之后，有必要把几何直觉的一切痕迹从实数定义中那个含糊的“量”的观念里清除出去。众所周知，这个问题大约就在那时几乎同时由康托尔、戴德金、梅雷和魏尔斯特拉斯用相当不同的方法解决了（见第~155~页）。
\end{enumerate}

从这时起，整数成了全部经典数学的基础。此外，随着公理方法的扩张和把数学对象视为心智的自由创造，基于算术的“模型”获得了更加重要的地位。对康托尔所宣称的这种自由，事实上仍然存在一个限制，那就是曾经萦绕于希腊人心头的“存在”问题；而它在这里以直接得多的方式出现，恰恰因为对直观表示的一切诉求此时都已被抛弃。我们将在后面（第~36~页以下）看到，“存在”概念在 20 世纪最初几年将成为怎样一个哲学-数学大漩涡的中心。但在 19 世纪，还没有走到那一步；要证明一个具有给定性质的数学对象存在，只需像欧几里得那样“构造”一个具有给定性质的对象。这恰恰正是算术“模型”的用处所在：一旦实数被“解释”成整数的说法，复数和欧氏几何也借助解析几何被如此解释，本世纪以来引入的所有新代数对象也都如此；最后——一项引起巨大轰动的发现——贝尔特拉米和克莱因甚至得到了罗巴切夫斯基和黎曼非欧几何的欧氏“模型”（见第~134~页），从而把这些乍看之下激起如此不信任的理论“算术化”了（并由此得到了完全的辩护）。

\begin{enumerate}
\def\labelenumi{\Alph{enumi})}
\setcounter{enumi}{3}
\tightlist
\item
  算术的公理化。——正是顺着这一演变，随后又转向了算术本身的根基，而事实上在大约 1880 年就可以看到这一点。看来在 19 世纪之前，除直接诉诸直觉之外，从未有人尝试定义整数的加法和乘法；莱布尼茨是唯一一个忠于自己原则、明确提醒人们 \(2 + 2 = 4\) 这样“自明”的真理只要想一想所涉各数的定义、便同样可以证明的人（{[}198 b{]}，第 IV 卷，第~403~页；参见 {[}69 a{]}，第~203~页）；而且他根本不认为加法和乘法的交换性是内禀的。\(^{36}\) 但他没有把这个想法再往前推进，而到 19 世纪中叶，这个方向上仍然毫无进展：就连其讲课大大助长了“算术化”观点传播的魏尔斯特拉斯，似乎也没有感到需要对整数理论作逻辑上的澄清。这个方向上的最初几步看来应归功于格拉斯曼，他在 1861 年（{[}134{]}，第 II 卷 \(_{2}\)，第~295~页）给出了整数加法和乘法的定义，并且只使用运算 \(x \mapsto x + 1\) 和归纳原理，证明了它们的基本性质（交换律、结合律、分配律）。归纳原理在 17 世纪已由 B. 帕斯卡清楚地构想并首次使用（{[}244{]}，第 III 卷，第~456~页）\(^{37}\)——尽管在古代也能找到它或多或少自觉的应用——而且从 17 世纪下半叶起就被数学家日常使用。但直到 1888 年，戴德金（{[}79{]}，第 III 卷，第~359-361~页）才陈述了算术的一个完备的公理系统（该系统三年后由皮亚诺重述，且通常冠以他的名字 {[}246 c{]}），其中特别包含归纳原理的一个精确表述（格拉斯曼使用了它，却没有明确陈述）。
\end{enumerate}

有了这种公理化，人们似乎已经达到了数学的最终基础。事实上，正当算术公理被清楚表述之时，对许多数学家（首先就是戴德金和皮亚诺本人）来说，这些公理就已经被剥夺了这门原始科学的角色，让位于数学诸理论中来得最晚的一个——集合论；而将要围绕整数概念展开的种种论争，是不能同 1900 年至 1930 年间那场大的“基础危机”分割开来的。
% semantic-fragments: legacy-input-seam
\relax{}
